\documentclass[10pt,a4paper]{amsart}
\usepackage[margin=19mm]{geometry}
\usepackage{amsmath,amssymb,mathtools}
\usepackage[scr=rsfs,cal=euler]{mathalfa}
\usepackage{xfrac}
\usepackage[unicode,hidelinks,bookmarks=false]{hyperref}
\newcommand{\ie}{i.e.\ }
\newcommand{\cf}{cf.\ }
\newcommand{\resp}{respectively\ }
\newcommand{\Subsection}{\subsection}
\providecommand{\nobreakdash}{\nobreak\hbox{-}}
\setlength{\emergencystretch}{2em}
\multlinegap0pt
% ---------------------------------------------------------------------------
% Editorial AMS wrapper prelude. The theorem environments and the mathematical macros below are
% transcribed from the paper's own preamble (cached-originals/source0.tex lines 74--320); the AMS amsart
% class, the named format packages of the pinned TeX Live runtime and this editorial formatting are not
% part of the author's text. No mathematics is changed.
% ---------------------------------------------------------------------------
\usepackage{dsfont}
\newtheorem{theorem}{Theorem}[section]
\newtheorem{lemma}[theorem]{Lemma}
\newtheorem{proposition}[theorem]{Proposition}
\newtheorem{assumption}[theorem]{Assumption}
\newtheorem{corollary}[theorem]{Corollary}
\newtheorem{definition}[theorem]{Definition}
\newtheorem{example}[theorem]{Example}
\newtheorem{remark}[theorem]{Remark}
\newtheorem{hypothesis}[theorem]{Hypothesis}
\newtheorem{claim}{Claim}
\renewcommand{\theequation}{\arabic{section}.\arabic{equation}}
\renewcommand{\thetheorem}{\arabic{section}.\arabic{theorem}}

\newcommand{\Tr}{\mathop{\mathrm{Tr}}}


\let\oldemptyset\emptyset
\let\emptyset\varnothing



\let\originalleft\left
\let\originalright\right
\renewcommand{\left}{\mathopen{}\mathclose\bgroup\originalleft}
\renewcommand{\right}{\aftergroup\egroup\originalright}


\theoremstyle{definition}

\newtheorem{condition}{Condition}[section]
\newtheorem{conjecture}{Conjecture}[section]
\newtheorem{criterion}{Criterion}[section]
\newtheorem{exercise}{Exercise}[section]
\newtheorem{notation}{Notation}[section]
\newtheorem{problem}{Problem}[section]
\newtheorem{solution}{Solution}[section]
\def\1{\Omega}
\def\t{t\wedge\tau_N^n}
\def\tt{T\wedge\tau_N^n}
\def\d{\mathrm{d}}
\def\dtau{\,\mathrm{d}\tau}
\def\dx{\,\mathrm{d}\boldsymbol{x}}
\def\dt{\,\mathrm{d}t}
\def\ds{\,\mathrm{d}s}
\def\dxt{\,\mathrm{d}\boldsymbol{x}\mathrm{d}t}
\def\dxs{\,\mathrm{d}\boldsymbol{x}\mathrm{d}s}
\def\dxtau{\,\mathrm{d}\boldsymbol{x}\mathrm{d}\tau}
\def\div{\operatorname{div}}
\def\divv{\mathbf{div}\,}
\def\supp{\operatorname{supp}}
\def\I{\mathrm{I}}
\def\B{\mathrm{B}}
\def\D{\mathrm{D}}
\def\A{\mathrm{A}}
\def\W{\mathrm{W}}
\def\R{\mathbb{R}}
\def\E{\mathrm{E}}
\def\Q{\mathrm{Q}}
\def\Z{\mathrm{Z}}
\def\loc{\mathrm{loc}}
\def\H{\mathbb{H}}
\def\X{\textsl{X}}
\def\V{\mathbb{V}}
\def\e{\varepsilon}
\def\L{\mathrm{L}}
\def\HH{\mathrm{H}}
\def\uu{\boldsymbol{u}}  % Gildo: Replaced \u by \vv in order to use the command \u already defined in latex for last i of Bogovskii name
\def\vv{\boldsymbol{v}}
\def\w{\boldsymbol{w}}
\def\n{\boldsymbol{n}}
\def\f{\boldsymbol{f}}
\def\00{\boldsymbol{0}}
\def\C{\mathrm{C}}
\def\F{\mathrm{F}}
\def\P{\mathbb{P}}
\def\N{\mathbb{N}}
\def\x{\boldsymbol{x}}
\def\y{\boldsymbol{y}}
\def\z{\boldsymbol{z}}
\def\G{\mathrm{G}}
\DeclareMathOperator*{\esssup}{ess\,sup}

\def\bfb{\mathbf{b}}
\def\bfB{\mathbf{B}}
\def\bfC{\mathbf{C}}
\def\bfX{\mathbf{X}}
\def\bfY{\mathbf{Y}}
\def\bfF{\mathbf{F}}
\def\bfG{\mathbf{G}}
\def\bfH{\mathbf{H}}
\def\bfL{\mathbf{L}}
\def\bfW{\mathbf{W}}
\def\bfQ{\mathbf{Q}}
\def\bfU{\mathbf{U}}
\def\bfS{\mathbf{S}}
\def\bfV{\mathbf{V}}
\def\bfI{\mathbf{I}}
\def\bfD{\mathbf{D}}
\def\bfQ{\mathbf{Q}}
\def\bfE{\mathbf{E}}
\def\bfG{\mathbf{G}}
\def\bfZ{\mathbf{Z}}
\def\bfe{\mathbf{e}}
\def\bfu{\mathbf{u}}
\def\bfv{\mathbf{v}}
\def\bff{\mathbf{f}}
\def\bfg{\mathbf{g}}
\def\bfa{\mathbf{a}}
\def\bfA{\mathbf{A}}
\def\bfE{\mathbf{E}}
\def\bfX{\mathbf{X}}
\def\bfb{\mathbf{b}}
\def\bfc{\mathbf{c}}
\def\bfh{\mathbf{h}}
\def\bf0{\mathbf{0}}
\def\bfz{\mathbf{z}}
\def\mcA{\mathcal{A}}
\def\bmcA{\boldsymbol{\mathcal{A}}}
\def\bmcB{\boldsymbol{\mathcal{B}}}
\def\bmcD{\boldsymbol{\mathcal{D}}}
\def\mcz{\mathcal{z}}

\def\rw{\rm{w}}


\def\bphi{\bm{\phi}}
\def\bPhi{\boldsymbol{\Phi}}
\def\bfi{\boldsymbol{\varphi}}
\def\bftau{\boldsymbol{\tau}}
\def\bpsi{\boldsymbol{\psi}}
\def\bxi{\boldsymbol{\xi}}
\def\bet{\boldsymbol{\eta}}
\def\bVcal{\boldsymbol{\mathcal{V}}}
\def\bMcal{\boldsymbol{\mathcal{M}}}

\def\Mcal{\mathcal{M}}

\def\bH{\boldsymbol{H}}
\def\bV{\boldsymbol{V}}
\def\bL{\boldsymbol{L}}
\def\bX{\boldsymbol{X}}
\def\bY{\boldsymbol{Y}}
\def\bW{\boldsymbol{W}}
\def\bbeta{\boldsymbol{\beta}}
\def\bp{\boldsymbol{p}}
\def\bh{\boldsymbol{h}}
\def\bB{\boldsymbol{B}}

\def\bD{\mathscr{D}}
\def\cD{\mathcal{D}}
\def\cB{\mathcal{B}}
\def\Var{\mathrm{Var}}

\DeclareMathAlphabet\mathbfcal{OMS}{cmsy}{b}{n}
\def\bcG{\mathbfcal{G}}
\def\bcH{\mathbfcal{H}}
\def\bcF{\mathbfcal{F}}

\def\Tr{\operatorname{Tr}}




\newcommand{\Addresses}{{% additional braces for segregating \footnotesize
		\footnote{
			\noindent \textsuperscript{1}Montanuniversit\"at Leoben, Department Mathematik und Informationstechnologie, Franz Josef Strasse 18, 8700
			Leoben, Austria.\\
			\textsuperscript{2}{FCT - Universidade do Algarve, Faro, Portugal \& CIDMA - Universidade de Aveiro, Aveiro, Portugal.}\\
			\textsuperscript{3}Department of Mathematics, Indian Institute of Technology Roorkee-IIT Roorkee,
			Haridwar Highway, Roorkee, Uttarakhand 247667, INDIA.			
			\par\nopagebreak
			\textit{e-mail:} \texttt{Ankit Kumar: ankitkumar.2608@gmail.com, ankit.kumar@unileoben.ac.at.}\\
				\textit{e-mail:} \texttt{Hermenegildo Borges de Oliveira: holivei@ualg.pt.}\\
			\noindent  \textit{e-mail:} \texttt{Manil T. Mohan: maniltmohan@ma.iitr.ac.in, maniltmohan@gmail.com.}
			
				$^\ast$Corresponding author.
			
			
			
			
			
			\textit{Keywords:} Generalised Navier-Stokes-Voigt equations, fluid dynamics, non-Newtonian fluids, weak solution.
			
			
			Mathematics Subject Classification (2020): {Primary 35Q35, 76A05, 76D03; Secondary 76D05, 35D30, 74F10.}
		
			
			
			
			
			
			
			
}}}
	

\begin{document}
\title{Optimal existence of weak solutions for the power-law Navier-Stokes-Voigt equations}
\author{Ankit Kumar, Hermenegildo Borges de Oliveira and Manil T. Mohan}
\date{Source: arXiv:2601.13051v3; distributed under CC BY 4.0}
\maketitle
\noindent\emph{Source, version and licence.} \emph{Optimal existence of weak solutions for the power-law Navier-Stokes-Voigt equations}, by Ankit Kumar, Hermenegildo Borges de Oliveira and Manil T. Mohan. The source of this editable unit is the e-print of \textbf{version v3} of arXiv:2601.13051, https://arxiv.org/abs/2601.13051v3, whose material is distributed under the Creative Commons Attribution 4.0 International licence, http://creativecommons.org/licenses/by/4.0/, as recorded in the arXiv licence metadata for this paper and shown on that abstract page. The whole of the body below is version v3, the newest version of the paper. Full rights notice: licenses/arxiv-2601.13051-CC-BY-4.0.txt.\par\smallskip
\noindent\textit{Editorial scope.} Counted unit: original Theorem 3.4 of the article, the statement carried by the source environment \texttt{theorem} with source label \texttt{UniT}, cached-originals/source0.tex lines 770--772, printed as \textbf{Theorem 3.4} exactly as in the article. It is delivered together with the authors' own complete proof as the single primary proof body, source0.tex lines 2499--2627 (129 lines): the \texttt{proof} environment whose optional argument names it "Proof of Theorem \textbackslash ref\{UniT\}", which carries out the whole uniqueness argument. No proof marker is added and no mathematics is written, repaired or re-derived; the authors' notation and the printed slips of the source are kept literal. No additional context is retained: the closure of the counted statement and proof over references and citations was computed and contains no reference and no citation at all, so the delivered passages are the statement and its complete proof, and every reference inside them resolves inside the delivered file. The article's own numbering is reproduced by explicit counter settings: the counter \texttt{theorem}, declared by \texttt{\textbackslash newtheorem\{theorem\}\{Theorem\}[section]} with lemma, proposition, assumption, corollary, definition, example, remark and hypothesis sharing it, is set before the retained passage, so the counted statement prints as Theorem 3.4 exactly as in the article. No \texttt{\textbackslash cite} command occurs in any retained passage, so this unit delivers no bibliography entry; the article's own bibliography data file \texttt{Power\_law\_NSV.bib} is neither spliced into the bundled source nor redistributed. The retained passages contain no figure environment and no graphics-inclusion command, so no artwork is redistributed. The source's own class is AMS \texttt{amsart} (\texttt{\textbackslash documentclass[12pt,reqno]\{amsart\}}); the e-print ships no class or style file, so no publisher class is shipped, used or redistributed, and the wrapper is the same AMS \texttt{amsart} class with the article's own macros and theorem declarations transcribed from its preamble. Every declared body of this unit is an exact, unmodified span of source0.tex: no body is a bounded passage comparison, \texttt{omit} was not used anywhere, and consequently no fidelity receipt is asserted in delivery-evidence.json. The complete concatenated original source is bundled as sources/192996/source0.tex. Omitted from this editable unit and therefore absent from it are source0.tex lines 1--338; 349--600; 604--675; 679--680; 701--712; 720--720; 735--736; 741--742; 751--769; 773--2498; 2628--2692, that is the article's own preamble, title and author block and abstract, the introduction, the two existence theorems and the whole approximation machinery of the intervening sections, the appendix, and the article's bibliography data, all of which remain present in the bundled source but are not delivered here. The AMS \texttt{amsart} wrapper, the named format packages of the pinned TeX Live runtime and this editorial source, licence and scope paragraph are editorial formatting only.\par\smallskip
\setcounter{section}{1}
\setcounter{equation}{0}
\setcounter{claim}{0}
\setcounter{condition}{0}
\setcounter{conjecture}{0}
\setcounter{criterion}{0}
\setcounter{exercise}{0}
\setcounter{notation}{0}
\setcounter{problem}{0}
\setcounter{solution}{0}
\setcounter{theorem}{0}
\begin{equation}\label{1.1}
	\left\{	
	\begin{aligned}
	{\partial_t(\vv -\kappa \Delta \vv ) +\divv(\vv\otimes \vv)+\nabla \pi -\nu\,\divv\!\left(|\bfD(\vv)|^{p-2}\bfD(\vv)\right)} & = {\f },\quad \text{ in }Q_T,  \\
		\div\vv &=0,\quad \text{ in }Q_T,  \\
		\vv &=\vv_0,\quad \text{ in }\1\times\{0\},  \\
		\vv &=\boldsymbol{0},\quad \text{ on }\Gamma_T,
	\end{aligned}
\right.
\end{equation}

\setcounter{section}{2}
\setcounter{equation}{5}
\setcounter{claim}{0}
\setcounter{condition}{0}
\setcounter{conjecture}{0}
\setcounter{criterion}{0}
\setcounter{exercise}{0}
\setcounter{notation}{0}
\setcounter{problem}{0}
\setcounter{solution}{0}
\setcounter{theorem}{1}
	\begin{align}\label{2.4}
		\big(|\bfD(\vv)|^{p-2}\bfD(\vv)-|\bfD(\uu)|^{p-2}\bfD(\uu)\big):\big(\bfD(\vv)-\bfD(\uu)\big)\geq 0,\  \text{ for all }\  \vv,\uu\in \bfV_p,
	\end{align}

\setcounter{section}{3}
\setcounter{equation}{0}
\setcounter{claim}{0}
\setcounter{condition}{0}
\setcounter{conjecture}{0}
\setcounter{criterion}{0}
\setcounter{exercise}{0}
\setcounter{notation}{0}
\setcounter{problem}{0}
\setcounter{solution}{0}
\setcounter{theorem}{0}
\begin{equation}\label{hyp:def:ws}
\f\in\bfL^1(Q_T),\qquad  \vv_0\in\bfV,\qquad p>1.
\end{equation}

\setcounter{section}{3}
\setcounter{equation}{1}
\setcounter{claim}{0}
\setcounter{condition}{0}
\setcounter{conjecture}{0}
\setcounter{criterion}{0}
\setcounter{exercise}{0}
\setcounter{notation}{0}
\setcounter{problem}{0}
\setcounter{solution}{0}
\setcounter{theorem}{0}
\begin{definition}\label{def.1}
{Assume that conditions \eqref{hyp:def:ws} are satisfied.}
We say that $\vv$ is a weak solution to the problem \eqref{1.1}, if:
\begin{enumerate}
  \item $\vv\in\L^\infty(0,T;\bfV)\cap\L^p(0,T;\bfV_p)$;
  \item ${\partial_t\vv}\in\L^2(0,T;\bfV)$;
  \item {$\vv(0)=\vv_0$;}
  \item The following identity,
\begin{equation}\label{3p3}
\begin{split}
& \hspace{-0.2cm}
-\!\int_{0}^{T}\!\!\!\!\int_{\Omega}\vv\cdot\partial_t\bfi \dxt - \kappa\!\int_{0}^{T}\!\!\!\!\int_{\Omega}\nabla\vv:\nabla\partial_t\bfi\dxt +
\nu\!\int_{0}^{T}\!\!\!\!\int_{\Omega}|\bfD(\vv)|^{p-2}\bfD(\vv):\bfD(\bfi)\dxt  \\
& \hspace{-0.2cm} =
\int_\1\vv_0\cdot\bfi(0)\dx + \kappa\int_\1\nabla\vv_0:\nabla\bfi(0)\dx + \!\int_{0}^{T}\!\!\!\!\int_{\Omega}\f\cdot\bfi\dxt + \!\int_{0}^{T}\!\!\!\!\int_{\Omega} \vv\otimes\vv:\mathds{\nabla}\bfi\dxt,
\end{split}
\end{equation}
holds for all $\bfi\in \bfC^\infty_{\div}(Q_T)$ with $\supp(\bfi) \Subset \Omega\times [0,T)$.
\end{enumerate}
\end{definition}

\setcounter{section}{3}
\setcounter{equation}{2}
\setcounter{claim}{0}
\setcounter{condition}{0}
\setcounter{conjecture}{0}
\setcounter{criterion}{0}
\setcounter{exercise}{0}
\setcounter{notation}{0}
\setcounter{problem}{0}
\setcounter{solution}{0}
\setcounter{theorem}{1}
\begin{alignat}{2}
\label{Hyp:u0}
  \vv_0&\in\bfV\cap\bfV_p,  \\
\label{Hyp:f:1}
 %\f\in\L^{\frac{2d}{d+2}}(0,T;\L^{\frac{2d}{d+2}}(\1)^d)
  \f&\in\L^{\infty}(0,T;\W^{-1,2}(\1)^d)\cap\L^{p'}(0,T;\L^{p'}(\1)^d).
\end{alignat}

\setcounter{section}{3}
\setcounter{equation}{4}
\setcounter{claim}{0}
\setcounter{condition}{0}
\setcounter{conjecture}{0}
\setcounter{criterion}{0}
\setcounter{exercise}{0}
\setcounter{notation}{0}
\setcounter{problem}{0}
\setcounter{solution}{0}
\setcounter{theorem}{1}
\begin{theorem}\label{thm:exist:1}
Let $\1\subset\R^d$ be a bounded domain with a \emph{Lipschitz-continuous boundary $\partial\1$}, and assume that hypotheses \eqref{Hyp:u0}--\eqref{Hyp:f:1} are verified.
Suppose that either
\begin{equation}\label{cond:1:thm1}
p>1\qquad \mbox{and}\qquad 2\leq d\leq 3,
\end{equation}
or
{
\begin{equation}\label{cond:2:thm1}
p\geq \frac{d}{2}\qquad \mbox{and}\qquad d>3.
\end{equation}
}
Then there exists, at least, a weak solution $\vv \in \C([0,T];\bfV)\cap \L^p(0,T;\bfV_p)$ to the problem \eqref{1.1}.
	\end{theorem}

\setcounter{section}{3}
\setcounter{equation}{6}
\setcounter{claim}{0}
\setcounter{condition}{0}
\setcounter{conjecture}{0}
\setcounter{criterion}{0}
\setcounter{exercise}{0}
\setcounter{notation}{0}
\setcounter{problem}{0}
\setcounter{solution}{0}
\setcounter{theorem}{2}
\begin{equation}\label{Hyp:f}
 %\f\in\L^{\frac{2d}{d+2}}(0,T;\L^{\frac{2d}{d+2}}(\1)^d)
  \f\in\L^{p'}(0,T;\L^{p'}(\1)^d).
\end{equation}

\setcounter{section}{3}
\setcounter{equation}{7}
\setcounter{claim}{0}
\setcounter{condition}{0}
\setcounter{conjecture}{0}
\setcounter{criterion}{0}
\setcounter{exercise}{0}
\setcounter{notation}{0}
\setcounter{problem}{0}
\setcounter{solution}{0}
\setcounter{theorem}{2}
\begin{theorem}\label{thm:exist:2}
Let $\1\subset\R^d$ be a bounded domain with a \emph{smooth boundary $\partial\1$ of class $\C^2$}, and assume that hypotheses \eqref{Hyp:u0} and \eqref{Hyp:f} are verified.
If
{\begin{equation}\label{3.7}
1<p<\frac{d}{2}\qquad \mbox{and}\qquad d>3,
\end{equation}
}then there exists, at least, a weak solution $\vv \in \C([0,T];\bfV)\cap \L^p(0,T;\bfV_p)$ to the problem \eqref{1.1}.
	\end{theorem}

\setcounter{section}{3}
\setcounter{equation}{8}
\setcounter{claim}{0}
\setcounter{condition}{0}
\setcounter{conjecture}{0}
\setcounter{criterion}{0}
\setcounter{exercise}{0}
\setcounter{notation}{0}
\setcounter{problem}{0}
\setcounter{solution}{0}
\setcounter{theorem}{3}
\begin{theorem}\label{UniT}
{	Under the assumptions of Theorems~\ref{thm:exist:1} and \ref{thm:exist:2}, for $p>1$ when $2\leq d\leq 4$, or for $p\geq \frac{d}{2}$ when $d>4$, the solution to the problem \eqref{1.1} is {unique in the class of weak solutions defined in Definition \ref{def.1}}.}
\end{theorem}

\setcounter{section}{7}
\setcounter{equation}{0}
\setcounter{claim}{2}
\setcounter{condition}{0}
\setcounter{conjecture}{0}
\setcounter{criterion}{0}
\setcounter{exercise}{0}
\setcounter{notation}{0}
\setcounter{problem}{0}
\setcounter{solution}{0}
\setcounter{theorem}{0}
\begin{proof}[Proof of Theorem~\ref{UniT}]
Let $\boldsymbol{v}_1, \boldsymbol{v}_2$ be two weak solutions  to the problem \eqref{1.1} in the sense of Definition~\ref{def.1}, with the same initial data $\boldsymbol{v}_0$ and forcing term $\boldsymbol{f}$.
Define
\begin{equation*}
\boldsymbol{w} := \boldsymbol{v}_1 - \boldsymbol{v}_2,
\end{equation*}
and note that $\div\w=0$ in $Q_T$ and $\w(0)=\bf0$ in $\Omega$.
Since $\w$ is not admissible as a test function in \eqref{3p3}, we introduce a mollification $\w^\varepsilon := \rho_\varepsilon * \w$, where the convolution is taken with respect to the time variable.
In particular, $\rho_\varepsilon(t)$ is supported in $(-\varepsilon,\varepsilon)$.
Then $\w^\varepsilon \in \bfC^\infty([0,T];\bfV_p)$, $\div\w^\varepsilon=0$, $\w^\varepsilon(\cdot,T)=0$ for $\varepsilon$ small enough, and
\begin{equation}
\label{sc:varep}
\w^\varepsilon \xrightarrow[\varepsilon\to0]{} \w,\quad \text{in}\quad \L^2(0,T;\bfV)\cap\L^p(0,T;\bfV_p).
\end{equation}
Using $\boldsymbol{w}^\varepsilon$ as a test function in \eqref{3p3} and subtracting the two weak formulations gives
\begin{align*}
& \int_0^t\!\!\int_\Omega \partial_\tau\boldsymbol{w} \cdot \boldsymbol{w}^\varepsilon \, \d \x\d\tau
+ \kappa \int_0^t\!\!\int_\Omega \nabla \partial_\tau\boldsymbol{w} : \nabla \boldsymbol{w}^\varepsilon \, \d \x\d\tau\ \\&\  \ +  \nu \int_0^t\!\!\int_\Omega \big(\bfA(\vv_1)-\bfA(\vv_2)\big):\mathbf{D}(\boldsymbol{w}^\varepsilon) \, \d \x\d\tau  \\
&= \int_0^t\!\!\int_\Omega (\boldsymbol{v}_1 \otimes \boldsymbol{v}_1 - \boldsymbol{v}_2 \otimes \boldsymbol{v}_2) : \nabla \boldsymbol{w}^\varepsilon \, \d \x\d\tau.
\end{align*}
Passing to the limit $\varepsilon \to 0$, using for this purpose \eqref{sc:varep}, this yields
\begin{equation*} %\label{dif:ws:v1-v2}
\begin{split}
& \frac{1}{2}\|\w(t)\|_{2}^2 + \frac{\kappa}{2}\|\nabla \w(t)\|_{2}^2 + \nu \int_0^t\!\!\int_\Omega \big(\bfA(\vv_1)-\bfA(\vv_2)\big):\bfD(\w)\,\d\x\d\tau  \\
&=  \int_0^t\!\!\int_\Omega (\vv_1\otimes\vv_1-\vv_2\otimes\vv_2):\mathds{\nabla}\w\d\x\d\tau.
\end{split}
\end{equation*}
\\
Combining the identity
\begin{equation*}
\vv_1\otimes\vv_1-\vv_2\otimes\vv_2= \vv_1\otimes\w + \w\otimes\vv_2
\end{equation*}
with the  monotonicity of $\bfA(\cdot)$ (see \eqref{2.4}), we obtain
\begin{equation}\label{dif:ws:v1-v2-1}
\begin{split}
& \frac{1}{2}\|\w(t)\|_{2}^2 + \frac{\kappa}{2}\|\nabla \w(t)\|_{2}^2 \leq
\int_0^t\!\!\int_\Omega \vv_1\otimes\w:\mathds{\nabla}\w\d\x\d\tau  + \int_0^t\!\!\int_\Omega \w\otimes\vv_2:\mathds{\nabla}\w\d\x\d\tau .
\end{split}
\end{equation}
Using integration by parts over $\Omega$, the facts that $\div\vv_1=0$ in $Q_T$ and $\vv_1=0$ on $\partial\Omega$ in the sense of traces, the first right-hand side term vanishes.
For the second right-hand side term, we divided the proof into two parts, which heavily rely on the dimensions, that $2\leq d \leq 4$ or $d>4$.

\vspace{2mm}
\noindent
\textbf{Part 1.} We first consider the case $p>1$ when $2\leq d  \leq 4$.
 For the second right-hand side term appearing in \eqref{dif:ws:v1-v2-1}, we use integration by parts, H\"older's and Sobolev's inequalities, and the fact that $\boldsymbol{v}_2\in\L^\infty(0,T;\bfV)$.
After all, we obtain
\begin{equation*}
\begin{split}
\left|\int_0^t \int_\Omega \boldsymbol{w} \otimes \boldsymbol{v}_2 : \nabla \boldsymbol{w}\d\x\d\tau\right| & =
\left|-\int_0^t \int_\Omega (\boldsymbol{w}\cdot\nabla)\boldsymbol{v}_2\cdot\boldsymbol{w}\d\x\d\tau \right|\leq
 \int_0^t \|\nabla\boldsymbol{v}_2(\tau) \|_2 \|\boldsymbol{w}(\tau) \|_4^2 \d\tau \\
& \leq  C\int_0^t \|\nabla\boldsymbol{w}(\tau) \|_2^2 \d\tau, \ \  d\leq 4,
\end{split}
\end{equation*}
for some positive constant $C$.
Using this information in \eqref{dif:ws:v1-v2-1}, we get
\begin{equation}\label{dif:ws:v1-v2:a}
\|\nabla \w(t)\|_{2}^2 \leq C\int_0^t \|\nabla\boldsymbol{w}(\tau) \|_2^2 \d\tau,
\end{equation}
for a distinct positive constant $C$ and $2\leq d\leq 4$.
Applying Gr\"onwall's inequality to \eqref{dif:ws:v1-v2:a}, we obtain
\begin{equation*}
	\|\nabla \w(t)\|_2^2 \leq \|\nabla \w(0)\|_2^2 e^{Ct},\ \text{ for all }\ \ t\in[0,T].
\end{equation*}
As $\w(0)=\boldsymbol{0}$, Sobolev's inequality implies that $\vv_1=\vv_2$ a.e. in $\Omega$ and for all $t\in [0,T]$ and hence the weak solution to the system \eqref{1.1} is unique for $p>1$ and $2\leq d\leq 4$.


\vspace{2mm}
\noindent
\textbf{Part 2.} Now, we consider the case $p\geq \frac{d}{2}$ when $d> 4$.
Again, we estimate the second right-hand side term appearing
in \eqref{dif:ws:v1-v2-1} using integration by parts, the fact
that $\div\w=0$ in $Q_T$, H\"older's and Sobolev's inequalities as
\begin{align}\label{P21}\nonumber
    & \left|\int_0^t \int_\Omega \boldsymbol{w} \otimes \boldsymbol{v}_2
    : \nabla \boldsymbol{w}\,\d\x\d\tau\right|
    =
    \left|-\int_0^t \int_\Omega (\boldsymbol{w}\cdot\nabla)
    \boldsymbol{v}_2\cdot\boldsymbol{w}\,\d\x\d\tau \right|
    \\& \leq  \nonumber
    \int_0^t \|\nabla\boldsymbol{v}_2(\tau) \|_{\frac{d}{2}}
    \|\boldsymbol{w}(\tau) \|_{\frac{2d}{d-2}}^2 \, \d\tau
  \leq
    C \int_0^t \|\nabla\boldsymbol{v}_2(\tau) \|_{\frac{d}{2}}
    \|\nabla \boldsymbol{w}(\tau) \|_{2}^2 \, \d\tau
    \\& \leq 
    C \int_0^t \|\nabla\boldsymbol{v}_2(\tau) \|_{p}
    \|\nabla \boldsymbol{w}(\tau) \|_{2}^2 \, \d\tau,
\end{align}
where in the last step we used the continuous embedding
$\L^p(\Omega)\hookrightarrow \L^{\frac{d}{2}}(\Omega)$,
which holds since $p\geq \frac{d}{2}$ and $\Omega$ is bounded.
Using \eqref{P21} in \eqref{dif:ws:v1-v2-1}, we find
\begin{equation}\label{P24}
    \|\nabla \w(t)\|_{2}^2 \leq
    C\int_0^t\|\nabla \boldsymbol{v}_2(\tau)\|_{p}
    \|\nabla\boldsymbol{w}(\tau) \|_2^2 \, \d\tau.
\end{equation}
Applying Gr\"onwall's inequality to \eqref{P24}, we get
\begin{align*}
    \|\nabla \boldsymbol{w}(t)\|_2^2
    &\leq \|\nabla \boldsymbol{w}(0)\|_2^2
    \exp \bigg\{ C\int_0^t \|\nabla \boldsymbol{v}_2(\tau)\|_{p}
    \, \d\tau \bigg\}.
\end{align*}
Since $\boldsymbol{v}_2 \in \L^p(0,T;\bfV_p)$ and
$p \geq \frac{d}{2} > 2$ for $d > 4$, H\"older's inequality
in time gives
\begin{align*}
    \int_0^T \|\nabla \boldsymbol{v}_2(\tau)\|_{p} \, \d\tau
    \leq T^{1 - \frac{1}{p}}
    \left( \int_0^T \|\nabla \boldsymbol{v}_2(\tau)\|_{p}^{p}
    \, \d\tau \right)^{\!\frac{1}{p}}
    < \infty.
\end{align*}
Therefore,
\begin{align*}
    \|\nabla \boldsymbol{w}(t)\|_2^2
    \leq \|\nabla \boldsymbol{w}(0)\|_2^2 \, e^{C},
    \quad \text{for all } t \in [0,T],
\end{align*}
for some finite constant $C = C(\nu, p, d, T, \boldsymbol{v}_2) > 0$.
As $\boldsymbol{w}(0)=\boldsymbol{0}$, Sobolev's inequality implies
$\boldsymbol{v}_1=\boldsymbol{v}_2$ a.e.\ in $\Omega$ and for all
$t\in [0,T]$.
Hence the weak solution to the system \eqref{1.1}
is unique for $p\geq \frac{d}{2}$ when $d>4$.
\end{proof}

\end{document}
